\documentclass[12pt]{exam}
\usepackage[utf8]{inputenc} 
\usepackage[brazil]{babel}
\usepackage{graphicx}
\usepackage{latexsym}
\usepackage{listings}
\usepackage{amsmath}
\usepackage{multirow}
\usepackage{multicol}
\usepackage{amsfonts}
\usepackage{amssymb}
\usepackage{epsfig,psfrag} 
\usepackage{wrapfig}
\usepackage[left=1.8cm,top=1cm,right=1.8cm,bottom=1cm]{geometry}
\usepackage[lined, ruled, vlined]{algorithm2e}
\usepackage[normalem]{ulem}
\usepackage{graphics}
\usepackage{pstricks,pst-node,pst-tree}
\usepackage{eqnarray,amsmath}
\usepackage{caption}
\usepackage{circuitikz}
\usepackage{datetime}
\usepackage{nicematrix, tikz}
\usepackage{mathtools}
\ctikzset{bipoles/thickness=1}
\usepackage[colorlinks=true, allcolors=blue]{hyperref}
\usepackage{listings}
\usepackage{xcolor}
\definecolor{cinza-codigo}{gray}{0.9}
\newcommand{\code}[1]{\colorbox{cinza-codigo}{\texttt{#1}}}
\usepackage{ragged2e}

% 2. Configurar o estilo do bloco de código
\lstset{
    language=bash,                  % Define a linguagem como Bash
    basicstyle=\ttfamily\small,     % Fonte monoespaçada e menor
    keywordstyle=\color{blue},      % Cor para comandos/palavras-chave
    commentstyle=\color{gray},      % Cor para comentários
    stringstyle=\color{purple},     % Cor para strings
    numbers=left,                   % Mostra números de linha na esquerda
    numberstyle=\tiny\color{gray},  % Estilo dos números de linha
    stepnumber=1,                   % Número em cada linha
    backgroundcolor=\color{lightgray!20}, % Cor de fundo do bloco
    frame=single,                   % Moldura ao redor do código
    breaklines=true,                % Quebra linhas longas automaticamente
    showstringspaces=false          % Não mostra espaços especiais em strings
}

\newcommand\encircle[1]{%
  \tikz[baseline=(X.base)] 
    \node (X) [draw, shape=circle, inner sep=0] {\strut #1};}

\newdateformat{specialdate}{\THEYEAR-\twodigit{\THEMONTH}-\twodigit{\THEDAY}}
\newcommand{\examdate}{Data de Entrega: 11/12/2024}

%\pagestyle{headandfoot}
%\firstpagefooter{}{Pág. \thepage\ de \numpages}{}
%\firstpageheader{}{}{}
%\runningfooter{}{Pág. \thepage\ de \numpages}{}
%\runningheader{}{}{}

\extraheadheight[.25in]{.5in}
\extrafootheight[.25in]{.5in}
\pagestyle{empty}
\usepackage{xcolor}
\definecolor{greendel}{HTML}{2F5416}
\definecolor{bluepoli}{HTML}{1E1D44}
\definecolor{redpoli}{HTML}{C73834}


\hypersetup{
    colorlinks=false,       % Desativa o texto colorido para manter a caixa de destaque
    linkbordercolor=green,  % Define a cor da caixa de links internos (sumário, referências) como verde
    urlbordercolor=red    % Define a cor da caixa de links externos (sites) como verde
}


% Definição de cores personalizadas
\definecolor{codegreen}{rgb}{0,0.6,0}
\definecolor{codegray}{rgb}{0.5,0.5,0.5}
\definecolor{codepurple}{rgb}{0.58,0,0.82}
\definecolor{backcolour}{rgb}{0.95,0.95,0.92}
\definecolor{codeblue}{rgb}{0,0,1}          % Define a cor azul pura

% Estilo customizado para Bash
\lstdefinestyle{bashcustom}{
    backgroundcolor=\color{cinza-codigo},   
    commentstyle=\color{codegreen},
    keywordstyle=\color{codeblue},
    %numberstyle=\tiny\color{codegray},
    stringstyle=\color{codeblue},
    basicstyle=\ttfamily\small,
    breakatwhitespace=false,         
    breaklines=true,                 
    captionpos=b,                    
    keepspaces=true,                 
    numbers=none,                    
    numbersep=5.5pt,                  
    showspaces=false,                
    showstringspaces=false,
    showtabs=false,                  
    tabsize=2,
    language=bash,
    morekeywords={cd, pwd, ls, mkdir, rm, cp, mv} 
}

\lstset{style=bashcustom}





\begin{document}
\setlength{\arrayrulewidth}{1.5pt}
\vspace{0.1cm}
\begin{tabular}{l|l}

\thispagestyle{empty}
\multirow{4}{30mm}{

    \vspace*{-1cm}\hspace*{-0.4cm}\includegraphics[width=0.15\textwidth]{logos/ciaa_gray.pdf}
}  
& {\bf \Large{Revisão Acelerada de Programação - RAP} \\
 & \bf \Large{RAP-2026, Marinha do Brasil} \\
 & {Centro de Instrução Almirante Alexandrino}\\
 & {Departamento de Aperfeiçoamento Avançado para Oficiais}\\
 & {Curso de Aperfeiçoamento Avançado em Guerra Acústica (C-ApA-Of-GA)}\\
 & {Professor: João Victor da Fonseca Pinto}\\
 & {Descrição do Documento: Configuração de ambiente VE no macOS}\\
 & {Código da Aula: 1}\\
%\hspace{0.0cm} & {Data: \today}\\
\end{tabular}\\ 
%\c

\thispagestyle{headandfoot}
\normalsize
\vspace{-0.2cm}
\vspace{0.2cm}



\section*{Configuração de ambiente}

No dia de hoje, realizaremos um avanço significativo no estabelecimento e na padronização de nosso ambiente prático de desenvolvimento de software. É frequente nos depararmos com cenários nos quais rotinas computacionais desenvolvidas em linguagem Python apresentam comportamento inconsistente ou falhas de execução ao serem transferidas entre diferentes estações de trabalho. Tais discrepâncias decorrem, em grande medida, de conflitos nas versões das bibliotecas instaladas ou da interferência do ambiente global do sistema operacional.

Para mitigar definitivamente essas inconsistências e assegurar a reprodutibilidade dos experimentos, adotaremos a combinação do Visual Studio Code (VS Code) com a ferramenta de isolamento virtualenv. O Visual Studio Code consolidou-se como um ambiente integrado de desenvolvimento (IDE) moderno, leve e altamente extensível, oferecendo suporte nativo à navegação de arquivos, integração com sistemas de controle de versão e execução de terminais de linha de comando em uma interface unificada.
Complementarmente, a utilização do virtualenv possibilita a criação de ambientes virtuais isolados dentro do próprio sistema, garantindo que o interpretador Python e suas respetivas dependências fiquem circunscritos ao escopo de cada projeto específico. Dessa forma, evita-se a contaminação do ambiente global e assegura-se que todo o corpo discente execute as atividades sob as exatas mesmas condições de configuração e versão.

\section*{Parte I}


Diferentemente do ambiente Windows, o sistema macOS não necessita de um subsistema Linux adicional (como o WSL), pois possui base UNIX nativa. Contudo, para disponibilizar utilitários essenciais de compilação e controle de versão no Terminal, faz-se necessária a instalação das \code{Xcode} e, em seguida, do gerenciador de pacotes \code{Homebrew}.

\subsection*{Instalação do Xcode Command Line Tools}

Abra o Terminal do macOS utilizando a busca do Spotlight (pressione as teclas \code{Cmd + Espaço}, digite \code{Terminal} e pressione \code{Enter}). No Terminal, execute o comando abaixo para dar início à verificação e instalação das ferramentas de linha de comando essenciais (que incluem o git e compiladores C/C++):
\vspace{0.3cm}
\begin{lstlisting}
xcode-select --install
\end{lstlisting}

Uma janela de diálogo do sistema será exibida na tela. Clique no botão \code{Instalar}, aceite os termos de licença e aguarde a conclusão do processo.

\subsection*{Instalação do Homebrew (Gerenciador de Pacotes):}

Caso a sua máquina não possua o gerenciador de pacotes \code{brew} instalado, execute o seguinte comando oficial no Terminal para baixá-lo e configurá-lo:

\vspace{0.3cm}
\begin{lstlisting}
/bin/bash -c "$(curl -fsSL https://raw.githubusercontent.com/Homebrew/install/HEAD/install.sh)"
\end{lstlisting}


Durante o processo, o instalador solicitará a senha de administrador da máquina. Digite-a (os caracteres não serão exibidos por questões de segurança) e pressione \code{Enter}. Ao término da instalação do Homebrew, preste atenção às instruções exibidas no próprio Terminal na seção \code{Next steps}. Em computadores com processadores da linha Apple Silicon, execute os comandos sugeridos para adicionar o brew à sua variável de ambiente PATH:

\vspace{0.3cm}

\begin{lstlisting}
echo 'eval "$(/opt/homebrew/bin/brew shellenv)"' >> ~/.zprofile
eval "$(/opt/homebrew/bin/brew shellenv)"
\end{lstlisting}


\section*{Parte II}

A escolha de um ambiente integrado dinâmico e flexível é fundamental para a edição de códigos, depuração e execução de scripts desenvolvidos para as atividades do curso. Para este propósito, adotaremos o \href{https://code.visualstudio.com}{Visual Studio Code (VS Code)}, um editor de código-fonte leve e extremamente extensível, mantido pela Microsoft.
Com o objetivo de viabilizar a conexão com ambientes isolados e garantir a perfeita integração com o interpretador da linguagem Python e servidores remotos, realizaremos a instalação do editor e a parametrização das extensões Python e Remote - SSH.

Para iniciar a instalação, acesse o portal eletrônico oficial do editor e efetue o \href{https://code.visualstudio.com/download?_exp_download=fb315fc982}{download} do instalador executável (User Installer) adequado à arquitetura do seu sistema. Efetue o download do arquivo de instalação correspondente ao macOS (Universal ou específico para o chip da sua máquina) e, ao concluir o download, arraste o aplicativo \code{Visual Studio Code.app} para a pasta Aplicativos (\code{Applications}) do seu sistema.


\section*{Parte III}

Com o propósito de viabilizar o desenvolvimento na linguagem Python e permitir a conexão com servidores remotos ou containers, faz-se necessária a adição de componentes de expansão ao editor. A interface gráfica do VS Code dispõe de um gerenciador nativo de extensões localizado na barra de ferramentas lateral esquerda, representado por um ícone com quatro blocos quadrados.


\subsection*{Configuração da Extensão Python}

Para habilitar o suporte completo à linguagem Python — incluindo destaque de sintaxe, autocompletar inteligente (IntelliSense) e recursos de depuração —, clique no ícone de Extensões na barra lateral esquerda ou utilize o atalho de teclado \code{Ctrl+Shift+X} (no Windows) ou \code{Cmd+Shift+X} (no macOS).
No campo de busca localizado no topo do painel lateral, digite exatamente o termo \code{Python}. O primeiro resultado listado corresponderá à extensão oficial desenvolvida pela Microsoft. Clique no botão \code{Install} associado ao pacote. O processo de transferência e ativação ocorrerá automaticamente em segundo plano. Após a conclusão, abra um arquivo de código com extensão \code{.py}.

\subsection*{Configuração da Extensão Remote - SSH}

Para efetuar a conexão com servidores remotos de processamento homologados ou máquinas virtuais através do protocolo SSH, utilizaremos o conjunto de extensões remotas.
Ainda no painel de Extensões, navegue até a caixa de pesquisa e digite \code{Remote - SSH}. Identifique o pacote oficial fornecido pela Microsoft e selecione o botão \code{Install}. Após o término da instalação, um novo ícone em formato de monitor ou conexão remota será adicionado à barra lateral do VS Code, representado graficamente por um pequeno botão azul contendo os símbolos de maior e menor justapostos (\code{><}).

\section*{Parte IV}

No macOS, utilizaremos o gerenciador de pacotes Homebrew para instalar a versão atualizada do Python 3, os utilitários de automação make, o controle de versão git e a ferramenta virtualenv. No Terminal do macOS, execute o seguinte comando:

\vspace{0.3cm}
\begin{lstlisting}
brew install python make git
pip3 install virtualenv
\end{lstlisting}

Aguarde a finalização da transferência e do desembalamento dos pacotes para garantir que o versionamento e o isolamento de dependências estejam devidamente operacionais.

\section*{Parte V}

Após a conclusão das instalações, navegue até o diretório de trabalho de sua preferência utilizando o terminal. Caso deseje criar uma nova pasta para organizar as atividades do curso, utilize o comando \code{mkdir} seguido pelo nome da pasta e acesse-a com o comando \code{cd}:

\vspace{0.3cm}
\begin{lstlisting}
mkdir -p ~/workspace
cd ~/workspace
\end{lstlisting}


Dentro do diretório escolhido, efetue a clonagem do repositório oficial da disciplina executando o seguinte comando no terminal:

\vspace{0.3cm}
\begin{lstlisting}
git clone https://github.com/jodafons/rap-2026.git
\end{lstlisting}


Concluído o download do projeto, acesse o diretório do repositório clonado e inicie o servidor do ambiente executando a instrução de automação via Make:

\vspace{0.3cm}
\begin{lstlisting}
cd rap-2026
make jupyter
\end{lstlisting}

Assim que a instrução for executada, a linha de comandos do terminal exibirá o log de inicialização do servidor, incluindo a URL de acesso local acompanhada do token de autenticação (por exemplo, \code{http://localhost:8888/lab}). Abra o navegador de internet de sua preferência (como Google Chrome, Mozilla Firefox ou Microsoft Edge) em sua máquina local, copie e cole o endereço de conexão fornecido pelo terminal na barra de navegação e pressione a tecla Enter. Este procedimento estabelecerá a conexão com a interface do Jupyter Notebook, viabilizando a interação com os arquivos interativos e notebooks de exercícios da aula.


\section*{Parte VI}

À medida que o curso avança, o professor disponibilizará novos notebooks, exercícios e atualizações no repositório remoto mantido no \href{https://github.com/jodafons/rap-2026/tree/main}{GitHub}. Para garantir que sua estação de trabalho local permaneça perfeitamente alinhada com a versão mais recente do material, disponibilizamos uma instrução simplificada de sincronização executada diretamente via \code{Make}.  

No entanto, faz-se estritamente necessário atentar para as regras de organização e preservação de arquivos no ambiente local. Todos os seus exercícios, testes, rotinas de aprendizado e anotações pessoais devem ser salvos obrigatoriamente dentro do diretório \code{rap-2026/notebooks/Aluno}. As informações e arquivos mantidos estritamente dentro dessa pasta não serão perdidos durante os procedimentos de sincronização. Por outro lado, ressalta-se que qualquer arquivo ou alteração realizada fora do diretório \code{notebooks/Aluno} será permanentemente perdido e sobrescrito no momento em que a atualização com o repositório do professor for efetuada.

Para realizar a atualização com segurança, abra o Terminal do macOS, navegue até o diretório do projeto com o comando \code{cd ~/workspace/rap-2026} e execute a instrução de automação \code{make update}. O comando \code{make update} efetuará o alinhamento de versões com o repositório remoto oficial, aplicando as melhorias e novos conteúdos da disciplina, ao mesmo tempo em que preserva intacto todo o trabalho armazenado na pasta \code{notebooks/Aluno}. Concluída a sincronização, reinicie o servidor do ambiente interativo executando a instrução \code{make jupyter}. Por fim, copie e cole o endereço de conexão fornecido no Terminal (como \code{http://localhost:8888/lab}) na barra de navegação do seu navegador para dar prosseguimento às atividades com o material devidamente atualizado.


\end{document}




